% !PRINTED 292
\weberpage{292}
Ist $\beta \equiv 0$, so hat (5) immer die eine Wurzel $z = \infty$; ist
dann $\delta \equiv \alpha$, und $\gamma$ nicht $\equiv 0$ $(\mathrm{mod}\ p)$, so gibt es nur diese
eine; in diesem Falle ist aber
\[
\alpha\delta \equiv 1, \qquad \alpha \equiv \delta, \qquad \left(\frac{\alpha+\delta}{2}\right)^2 \equiv 1 \quad (\mathrm{mod}\ p),
\]
und die Bedingung I. erfüllt. Ist aber gleichzeitig $\gamma \equiv 0$ $(\mathrm{mod}\ p)$,
so ist die Substitution die \gesperrt{identische}.

Ist $\beta \equiv 0$, aber $\alpha-\delta$ nicht $\equiv 0$ $(\mathrm{mod}\ p)$, so hat (5) noch
eine zweite Wurzel; da jetzt $\alpha\delta \equiv 1$ $(\mathrm{mod}\ p)$, so ist:
\[
\left(\frac{\alpha+\delta}{2}\right)^2 - 1 \equiv \left(\frac{\alpha-\delta}{2}\right)^2 \quad (\mathrm{mod}\ p),
\]
also quadratischer Rest, und die Bedingung II. erfüllt. Wir
fassen diese Sätze so zusammen:

\gesperrt{Im Falle I. bleibt ein Element oder alle Elemente
ungeändert, im Falle II. bleiben zwei Elemente ungeändert
 und im Falle III. werden alle Elemente
geändert.}

Wenn wir eine und dieselbe Substitution
\weq{7}{A = \begin{pmatrix}\alpha, & \beta \\ \gamma, & \delta\end{pmatrix}}
mehrmals wiederholen, so entstehen die Substitutionen
\[
A, \ A^2, \ A^3 \dots,
\]
in deren Reihe einmal die identische Substitution $\begin{pmatrix}1, & 0 \\ 0, & 1\end{pmatrix}$
auftreten muß. Ist $n$ die kleinste positive Zahl, für die
\[
A^n = \begin{pmatrix}1, & 0 \\ 0, & 1\end{pmatrix}
\]
ist, so heißt $n$ der Grad von $A$ (Bd. II, \S\,2).

Setzen wir für ein beliebiges $m$
\[
A^m = \begin{pmatrix}\alpha_m, & \beta_m \\ \gamma_m, & \delta_m\end{pmatrix},
\]
so erhalten wir zur Berechnung der Zahlen $\alpha_m$, $\beta_m$, $\gamma_m$,
$\delta_m$ folgendes System rekurrenter Formeln:
\weq{8}{\begin{aligned}
\alpha_{m+1} &= \alpha\alpha_m+\beta\gamma_m, & \beta_{m+1} &= \alpha\beta_m+\beta\delta_m,\\
\gamma'_{m+1} &= \gamma\alpha_m+\delta\gamma_m, & \delta_{m+1} &= \gamma\beta_m+\delta\delta_m.
\end{aligned}} % [sic] printed with a prime on \gamma'_{m+1}, inconsistent with the other three terms

Besonders einfach lassen sich hieraus die Zahlen $\alpha_m$, $\beta_m$,
$\gamma_m$, $\delta_m$ im Falle I. berechnen.
